\begin{ouverture}[courbe]{5}{Équations, inéquations et produits nuls}{Partie 1 · Nombres et calculs}
{Équations équivalentes · Inéquations · Produits nuls · Analyse-synthèse}
\objectif{Résoudre une équation ou une inéquation du premier degré.}
\objectif{Factoriser pour résoudre une équation produit nul.}
\objectif{Mettre un problème en équation et valider sa solution.}
\prerequis{Calcul littéral, identités remarquables et opérations inverses.}
\end{ouverture}
\pagestyle{otMain}\essentiel\fontsize{11}{15}\selectfont
\essSection{1}{Transformer une équation}
\begin{definition}Résoudre une équation, c’est trouver toutes les valeurs de l’inconnue qui rendent l’égalité vraie. Deux équations sont équivalentes lorsqu’elles ont les mêmes solutions.\end{definition}
\begin{propriete}Ajouter ou soustraire une même quantité aux deux membres donne une équation équivalente. Multiplier ou diviser les deux membres par un même nombre non nul aussi.\end{propriete}
\begin{exemple}$5x-7=2x+8$ équivaut à $3x=15$, donc à $x=5$. Contrôle : les deux membres initiaux valent 18.\end{exemple}
\begin{savoirfaire}[Analyse et synthèse]Analyse : on suppose qu’un nombre est solution et on en déduit les seules valeurs possibles. Synthèse : on vérifie que ces valeurs conviennent, dans l’équation et dans le contexte. Cette vérification est indispensable lorsque toutes les étapes ne sont pas des équivalences.\end{savoirfaire}
\begin{attention}Une équation peut avoir une seule solution, aucune solution, ou tous les nombres pour solutions. $2x+3=2x+5$ est impossible ; $2(x+3)=2x+6$ est toujours vraie. Diviser par $x$ peut perdre la solution $x=0$.\end{attention}
\clearpage
\essSection{2}{Inéquations et produits nuls}
\begin{propriete}Ajouter un même nombre aux deux membres d’une inégalité conserve son sens. Multiplier ou diviser par un nombre strictement positif le conserve aussi ; par un nombre strictement négatif, il change le sens.\end{propriete}
\begin{exemple}$-2x\ge8$ équivaut à $x\le-4$, car on divise par $-2$. Sur une droite graduée, les solutions forment la demi-droite à gauche de $-4$, le crochet tourné vers les solutions indiquant que $-4$ est inclus. On peut noter cet ensemble $]-\infty\,;-4]$.\end{exemple}
\begin{savoirfaire}[Résoudre graphiquement $ax\ge b$]Tracer la droite d’équation $y=ax$ et la droite horizontale $y=b$. Les solutions sont les abscisses des points de la première droite situés sur ou au-dessus de la seconde.\end{savoirfaire}
\begin{propriete}Un produit est nul si et seulement si l’un au moins de ses facteurs est nul. Ainsi $(x-4)(2x+3)=0$ donne $x=4$ ou $x=-3/2$.\end{propriete}
\begin{savoirfaire}[Se ramener à un produit nul]Tout passer dans un membre, puis factoriser : $x^2=5x$ équivaut à $x(x-5)=0$ ; $(x+1)^2-9=0$ équivaut à $(x-2)(x+4)=0$ grâce à $a^2-b^2$.\end{savoirfaire}
\clearpage\rubrique{Découverte}
\begin{blocLivre}{ex-5-1}\begin{exercice}[1]{Une opération inverse}
Résoudre $x+13=4$ et $-4x=18$.
\end{exercice}\end{blocLivre}
\begin{blocLivre}{ex-5-2}\begin{exercice}[1]{Deux étapes}
Résoudre $7x-5=30$.
\end{exercice}\end{blocLivre}
\begin{blocLivre}{ex-5-3}\begin{exercice}[1]{L’inconnue des deux côtés}
Résoudre $8x+3=5x-12$.
\end{exercice}\end{blocLivre}
\begin{blocLivre}{ex-5-4}\begin{exercice}[1]{Avec une parenthèse}
Résoudre $3(2x-5)=4x+9$.
\end{exercice}\end{blocLivre}
\begin{blocLivre}{ex-5-5}\begin{exercice}[1]{Une première inéquation}
Résoudre $3x\ge12$, puis $-2x\ge8$. Représenter chaque ensemble de solutions sur une droite graduée.
\label{fig-ex-5-5}\figureGeo{inequation-droites}
\end{exercice}\end{blocLivre}

\clearpage\rubrique{Entraînement}
\begin{blocLivre}{ex-5-6}\begin{exercice}[2]{Une équation fractionnaire}
Résoudre $\dfrac{x}{3}+\dfrac{x}{4}=14$.
\end{exercice}\end{blocLivre}
\begin{blocLivre}{ex-5-7}\begin{exercice}[2]{Tout ou rien}
Résoudre $4(x-2)=4x-8$, puis $4(x-2)=4x+1$.
\end{exercice}\end{blocLivre}
\begin{blocLivre}{ex-5-8}\begin{exercice}[2]{Un produit nul}
Résoudre $(2x-7)(x+4)=0$.
\end{exercice}\end{blocLivre}
\begin{blocLivre}{ex-5-9}\begin{exercice}[2]{Factoriser pour résoudre}
Résoudre $x^2-6x=0$, puis $(x+3)^2-16=0$.
\end{exercice}\end{blocLivre}
\begin{blocLivre}{ex-5-10}\begin{exercice}[2]{Lire une inéquation sur un graphique}
La droite représente la fonction $x\mapsto1{,}5x$. Résoudre graphiquement $1{,}5x\ge6$, puis confirmer par le calcul.
\label{fig-ex-5-10}\figureGeo{inequation-graphique}
\end{exercice}\end{blocLivre}

\clearpage\rubrique{Problèmes}
\begin{blocLivre}{ex-5-11}\begin{exercice}[3]{Des billets de spectacle}
Un groupe achète 28 billets : les billets jeunes coûtent 6 €, les autres 10 €. La dépense totale est 232 €. Combien de billets jeunes ?
\end{exercice}\end{blocLivre}
\begin{blocLivre}{ex-5-12}\begin{exercice}[3]{Un rectangle à retrouver}
Un rectangle a un périmètre de 58 cm. Sa longueur dépasse sa largeur de 7 cm. Déterminer ses dimensions et son aire.
\label{fig-ex-5-12}\figureGeo{rectangle-equation}
\end{exercice}\end{blocLivre}
\begin{blocLivre}{ex-5-13}\begin{exercice}[3]{Un achat rentable}
Une gourde filtrante coûte 38 € et ses recharges 12 € par an. Elle remplace un achat de bouteilles estimé à 30 € par an. Écrire une inéquation traduisant « après $n$ années, la dépense avec la gourde est inférieure ou égale à celle des bouteilles », puis la résoudre.
\end{exercice}\end{blocLivre}
\begin{blocLivre}{ex-5-14}\begin{exercice}[3]{Deux aires égales}
Un carré de côté $x+2$ cm et un rectangle de côtés $x$ cm et $x+5$ cm ont la même aire. Les dimensions sont positives. Déterminer $x$ et l’aire commune.
\label{fig-ex-5-14}\figureGeo{aires-egales}
\end{exercice}\end{blocLivre}
\begin{blocLivre}{ex-5-15}\begin{exercice}[3]{Analyse et synthèse}
Un rectangle a une longueur qui dépasse sa largeur $x$ de 2 cm et une aire de 24 cm$^2$. Vérifier que $x(x+2)-24=(x+1)^2-25$, puis déterminer ses dimensions par une analyse suivie d’une synthèse.
\end{exercice}\end{blocLivre}

\clearpage
